\documentclass{handout}

\assignment{Homework 1}
\title{Motion in a straight line}
% The original date was incomplete — it read "Due Monday,  2025" with no month
% or day. Fill in and uncomment.
% \duedate{Monday, September 8, 2025}

\begin{document}
\maketitle

\hwinstructions

\begin{questions}
    \question \textbf{Average walking velocity:} Estimate your average velocity when walking from the dorms to Straz. Explain how you measured the distance and the time it takes, also explain how you used these quantities to calculate your average walking velocity. Express your answer both in meters per second and feet per second.
    \question \textbf{Derivative product rule:} In class we learned how to obtain derivatives of polynomials: $$\frac{d}{dt} t^n=nt^{n-1}.$$
    We did this mostly with drawings. But what is the derivative of two functions multiplied together, i.e. $f(t)=g(t)h(t)$? An example would be $f(t)=t^2 *\sin(t)$. Use a rectangle of side lengths $g$ and $h$ and similar reasoning from class to show that
    $$
    \frac{d}{dt}f(t)=\left(\frac{d}{dt}g(t)\right)h(t) + g(t)\left(\frac{d}{dt}h(t)\right).
    $$
    \question \textbf{Building height:} Physics lets use a stopwatch as a measuring stick to measure tall buildings. You can use ft/s ($g$=32 ft/s$^2$) or m/s ($g$=9.8m/s$^2$). Say you drop a ball and it takes 2 s to hit the ground.
    \begin{parts}
        \part Draw a velocity-time graph for a ball dropped off the top of a building (\textit{hint}: what is the initial velocity?).
        \part What is the velocity of the ball when it hits the ground?
        \part Draw a position-time graph for the ball, let your position on top of the building be $z=0$ (\textit{hint}: what is the initial position and how does the initial velocity affect the graph?).
        \part Calculate the height of the building.
        \part If you go to a different building where the ball takes twice as long to reach the ground, how much taller is the building?
    \end{parts}

    \question \textbf{Height of ball thrown straight up:} The analysis for a ball thrown straight up is a little different. All that is known is that the initial and final \textit{speed} are the same (their velocities are opposite). Suppose the time from throwing to catching is again 2 s.
    \begin{parts}
        \part Draw a velocity-time graph (the initial velocity is unknown but positive).
        \part Use your diagram and the total time to find the initial velocity.
        \part Use the initial velocity and the time to maximum to find out how high the ball went.
    \end{parts}
\end{questions}

\end{document}
